% Begin


%%%%%%%%%%%%%%%%%%%%%%%%%%%%%%%%%%%%%%%%%%%%%%%%%%%%%%%%%%%%%%%
\chapter*{\S \space 31. --- DIGRESSION SUR LES RELÈVEMENTS D'UNE ACTION EXTÉRIEURE D'UN GROUPE FINI $G$ SUR UN GROUPE PROFINI À LACETS $\pi$}\thispagestyle{empty}
\addcontentsline{toc}{section}{31. Digression sur les relèvements d'une action extérieure d'un groupe fini sur un groupe profini à lacets}
\label{sec:31}
\section*{}

On suppose l'action extérieure \emph{fidèle}, i.e. $G \subset \hat{\hat{\mathfrak{T}}} (\pi)=\hat{\hat{\mathfrak{T}}}$, et que $G = G^+$.  On suppose de plus qu'il existe une discrétification invariante $\pi_0$, i.e. telle que $G\subset \mathfrak{T}(\pi_0)$.  Si on suppose $G$ cyclique, alors sauf erreur il est prouvé que la situation est \emph{réalisable} topologiquement (donc aussi de fa\c{c}on analytique complexe\dots)

Dans le cas discret, la signification des classes de $\pi_0$ conjugaison de relèvement $G \to \mathfrak{S}(\pi_0)$ est bien comprise, de même que pour les relèvements partiels. On trouve une réalisation canonique [si $U^!\ne\emptyset$] du groupoïde fondamental associé au groupe extérieur $\pi_0$, par un groupoïde fini $[U^!]$ sur lequel $G$ opère au sens strict, dont les points correspondent aux classes de $\pi_0$-conjugaison de sections partielles $\ne 1$ maximales.  Le groupe $\mathfrak{T}(\pi_0)^G$ opère de fa\c con également canonique sur ce groupoïde. Si $P \in U^!$ correspond à un relèvement $G\to\mathfrak{S}(\pi_0)$, alors un élément $\cdot u\in\mathfrak{T}(\pi_0)^G$ est dans l'image de $\mathfrak{S}(\pi_0)^G$, i.e. peut se remonter en $u\in\mathfrak{S}(\pi_0)$ commutant à $G$, si et seulement si $\cdot u(P)=P$, i.e. (si $\cdot u$ est remonté de fa\c{c}on quelconque en $v$), si et seulement si $v(G)$ est $\pi_0$-conjugué à $G$, i.e. si et seulement s'il existe $g\in\pi_0$ tel que $v(G)= \operatorname{int}(g)(G)$, i.e. $u \defeq \operatorname{int}(g^{-1})v$ fixe $G$ (auquel cas bien sûr il centralise, par l'hypothèse $\cdot u \in \mathfrak{T}^G$). Donc notre brillante assertion est une tautologie, et donc le relèvement est unique. Comme $U^!$ est fini, le stabilisateur $\mathfrak{T}(\pi_0)^G_P$ de $P$ dans $\mathfrak{T}(\pi_0)^G$ est d'indice fini, et c'est donc ce sous-groupe d'indice fini qui se remonte gaillardement et canoniquement.

Que peut-on dire dans le contexte profini~?  Bien sûr, on a une application canonique
\vskip .2cm
classes de $\pi_0$-conjugaison\ \ \ \ \ \ \ \ \ \ \ \ \ \ \ \ 
\ \ \ \ \ \ \  classes de $\pi_0$-conjugaison

de relèvements de \ \ \ \ \ \ \ \ \ \ \ \ \ \ \ \ \ \ \ \ \ \ \ \ \ \ \ \ \ 
\ \ \ des relèvements de

$G$ en $G\to\mathfrak{S}(\pi_0)$\ \ \ \ \ \ \ \ \ \ \ \ \ \ \  \ \ \ \
$\longrightarrow$
\ \ \ \ \ \ \ \ \ \ $G$ en $G\to \hat{\hat{\mathfrak{S}}}(\pi_0)$,

i.e. de scindage de l'extension\ \ \ \ \ \ \ \ \ \ \ \ \ \ \ \ \ \ i.e. de scindage de l'extension

$1\to\pi_0\to E\to G \to 1$\ \ \ \ 
\ \ \ \ \ \ \ \ \ \ \ \ \ \ \ \ \ \ \ \ \ \ \ \ $1\to\hat\pi_0=\pi
\to \hat{E} \to G \to 1$.
\vskip .2cm
Il faudrait examiner d'abord
\begin{enumerate}
    \item[a)] la question de la bijectivité de cette application,
    \item[b)] si $\pi^G=1$ (rigidité), pour $G\ne \{1\}$. (pour un relèvement donné, dans $E$ pour simplifier).
\end{enumerate}
J'ai envie de conjecturer sans vergogne qu'il en est ainsi --- ce qui impliquerait par exemple que pour tout tel relèvement de $G$, il y a un sous-groupe \emph{ouvert} $(\mathfrak{T}^G)_P$ du groupe (peut-être vraiment immense a priori~!) $\hat{\hat{\mathfrak{T}}}^G$, qui se remonte canoniquement de fa\c{c}on à commuter à $G$\dots.

Il faudrait manifestement faire, en même temps qu'une théorie des opérations extérieures des groupes finis sur des groupes discrets à lacets (qui est pour le moment extrêmement conjecturale) une théorie analogue dans le cas profini --- j'aurai sans doute à y revenir par la suite.


% End
